\documentclass{article}
\usepackage[utf8]{inputenc}
\usepackage{amsmath}
\usepackage{geometry}
\geometry{margin=2cm}
\begin{document}

\section*{Análise Consolidada da Questão 145 — ENEM 2024 — Matemática}

\subsection*{Contextualização}
O campo do estádio Maracanã foi adaptado para atender aos padrões da FIFA para a Copa do Mundo de 2014, reduzindo-se suas dimensões de \(110\text{ m} \times 75\text{ m}\) para \(105\text{ m} \times 68\text{ m}\). A questão propõe calcular a redução da área do campo resultante dessas mudanças.

\section{Solução Técnica}

\subsection{Resposta Final}
A resposta correta é \textbf{1110} metros quadrados.

\subsection{Estratégia}
Calcular a área original do campo e a nova área após a redução das dimensões e depois encontrar a diferença entre essas áreas.

\subsection{Passos para a solução}
\begin{enumerate}
  \item Calcular a área original do campo com dimensões \(110\text{ m} \times 75\text{ m}\):
  \[
    A_{\text{original}} = 110 \times 75 = 8250\text{ m}^2.
  \]
  \item Calcular a nova área do campo com dimensões \(105\text{ m} \times 68\text{ m}\):
  \[
    A_{\text{novo}} = 105 \times 68 = 7140\text{ m}^2.
  \]
  \item Calcular a redução na área:
  \[
    \Delta A = 8250 - 7140 = 1110\text{ m}^2.
  \]
\end{enumerate}

\subsection{Fórmulas Utilizadas}
\begin{equation}
\text{Área do retângulo} = \text{comprimento} \times \text{largura}.
\end{equation}

\subsection{Dificuldades comuns e dicas didáticas}
Para ajudar alunos, é importante destacar que a redução de área não é dada pela diferença simples das dimensões, mas pela diferença dos produtos dessas dimensões. Reforçar o conceito da área do retângulo e fazer exercícios com números menores pode facilitar a compreensão.

\section{Análise dos Distratores}

\begin{itemize}
  \item \textbf{Opção A}: Resultado muito baixo devido a erro de cálculo ou interpretação errada das dimensões.
  \item \textbf{Opção B}: Confunde diferença linear com diferença de área.
  \item \textbf{Opção C}: Subtrai uma dimensão ou fração incorretamente.
  \item \textbf{Opção E}: Erro complexo mas próximo, resultado de cálculo incorreto.
\end{itemize}

\section{Perfil da Questão}
A questão exige habilidades básicas em cálculo de área de retângulos, interpretação quantitativa do enunciado e atenção às unidades.

\section{Conclusão de Consistência}
A análise confirmou que a solução está correta e coerente, o raciocínio é consistente e os distractores são plausíveis. A questão é adequada para o nível definido e não apresenta incoerências.

\end{document}